\setcounter{section}{21}

  \zwischenueberschrift{Übungsaufgaben}

\inputaufgabe
{}
{

Beschreibe diverse Kleidungsstücke als zweidimensionale
\definitionsverweis {Mannigfaltigkeiten mit Rand}{}{.}

}
{} {}

\inputaufgabe
{}
{

Definiere die Konzepte differenzierbare Abbildung, Diffeomorphismus, Tangentialraum, Tangentialbündel ... für
\definitionsverweis {Mannigfaltigkeiten mit Rand}{}{.}

}
{} {}

\inputaufgabe
{}
{

Beschreibe
\definitionsverweis {reelle Intervalle}{}{}
als
\definitionsverweis {Mannigfaltigkeiten mit Rand}{}{.}
Was ist jeweils der
\definitionsverweis {Rand}{}{,}
welche sind untereinander
\definitionsverweis {diffeomorph}{}{.}

}
{} {}

\inputaufgabe
{}
{

$\,$

\bild{ \begin{center}
\includegraphics[width=5.5cm]{ \bildeinlesungsvg {Circle_on_sphere_wireframe_10deg_6r} {svg} }
\end{center}
\bildtext {} }

\bildlizenz { Circle on sphere wireframe 10deg 6r.svg } {} {Itai} {Commons} {CC-by-sa 3.0} {}

Zeige, dass sowohl das blaue als auch das rote Oberflächenstück einschließlich der Begrenzungslinie eine
\definitionsverweis {Mannigfaltigkeit mit Rand}{}{}
ist. Was ist der Rand? Sind die beiden Mannigfaltigkeiten
\definitionsverweis {diffeomorph}{}{?}
Gibt es eine einfachere dazu diffeomorphe Mannigfaltigkeit?

}
{} {}

\inputaufgabe
{}
{

Es sei $M$ eine
\definitionsverweis {differenzierbare Mannigfaltigkeit}{}{}
\zusatzklammer {ohne Rand} {} {}
und $N$ eine
\definitionsverweis {differenzierbare Mannigfaltigkeit mit Rand}{}{.}
Was kann man über das
\definitionsverweis {Produkt}{}{}
\mathl{M \times N}{} sagen?

}
{} {}

\inputaufgabe
{}
{

Zeige, dass der abgeschlossene Annulus

\mathdisp {\left\{ P \in \R^2  \mid   1 \leq \Vert { P } \Vert \leq 2         \right\}} {  }

und der abgeschlossene Zylinder

\mathdisp {S^1 \times [0,1]} {  }

zueinander
\definitionsverweis {diffeomorphe}{}{} \definitionsverweis {Mannigfaltigkeiten mit Rand}{}{}
sind.

}
{} {}

\inputaufgabegibtloesung
{}
{

Begründe, dass der $\R^2$ keine Struktur einer
\definitionsverweis {Mannigfaltigkeit mit Rand}{}{}
derart trägt, dass die angegebene Teilmenge $T$ der
\definitionsverweis {Rand}{}{}
ist.
\aufzaehlungvierabc{
\mavergleichskette
{\vergleichskette
{ T
}
{ = }{ {]0,1[}
}
{  }{
}
{    }{
}
{   }{
}
}
{}{}{,}
wobei das Intervall auf der $x$-Achse liegt.
}{
\mavergleichskette
{\vergleichskette
{ T
}
{ = }{ [0,1]
}
{  }{
}
{    }{
}
{   }{
}
}
{}{}{,}
wobei das Intervall auf der $x$-Achse liegt.
}{
\mavergleichskette
{\vergleichskette
{ T
}
{ = }{ \R
}
{  }{
}
{    }{
}
{   }{
}
}
{}{}{,}
wobei $\R$ die $x$-Achse sei.
}{
\mavergleichskette
{\vergleichskette
{ T
}
{ = }{ S^1
}
{  }{
}
{    }{
}
{   }{
}
}
{}{}{.}
}

}
{} {}

\inputaufgabe
{}
{

Zeige, dass die zueinander inversen
\definitionsverweis {Diffeomorphismen}{}{}
\maabbeledisp {\varphi} { {\mathbb H} } { U \left(   0 , 1  \right)
} { z} { { \frac{  z-  { \mathrm i}  }{ z + { \mathrm i}  } }
} {,}
und
\maabbeledisp {\psi} { U \left(   0 , 1  \right)  } { {\mathbb H}
} { w } { { \frac{   (w+1)  { \mathrm i}  }{  1-w  } }
} {,}
\zusatzklammer {siehe
[[Obere Halbebene/Einheitskreis/Rational isomorph/Fakt/Beweis/Aufgabe|Aufgabe 18.10]]} {} {}
zu Diffeomorphismen
\zusatzklammer {zwischen
\definitionsverweis {Mannigfaltigkeiten mit Rand}{}{}} {} {}
zwischen der abgeschlossenen Halbebene und der abgeschlossenen punktierten Kreisscheibe
\mathl{B  \left(  0 , 1 \right) \setminus \{  \left(  1   , \,   0                   \right) \}}{} ausgedehnt werden können.

}
{} {}

\inputaufgabe
{}
{

Es sei $M$ eine
\definitionsverweis {differenzierbare Mannigfaltigkeit mit Rand}{}{.}
Unter einem \stichwort {differenzierbaren Halbweg} {} verstehen wir jede
\definitionsverweis {differenzierbare Abbildung}{}{}
\maabbdisp {\gamma} { [0, \epsilon [ } { M
} {}
oder
\maabbdisp {\gamma} {] -\epsilon, 0] } { M
} {}
\zusatzklammer {mit \mathlk{\epsilon >0}{.} Sie heißen nach innen bzw. nach außen gerichtet} {} {.}
Definiere, wann zwei Halbwege mit
\mathl{\gamma_1(0) = \gamma_2(0) =P \in M}{} \stichwort {tangential äquivalent} {} sind, und zeige, dass dadurch eine Äquivalenzrelation gegeben ist, und dass die
\definitionsverweis {Quotientenmenge}{}{}
ein reeller Vektorraum ist, der der \stichwort {Tangentialraum} {} in $P$ heißt. Charakterisiere die Äquivalenzklassen, die sowohl nach innen als auch nach außen repräsentierbar sind.

}
{} {}

\inputaufgabegibtloesung
{}
{

Es sei $M$ eine
\definitionsverweis {Mannigfaltigkeit mit Rand}{}{}
und sei

\mavergleichskette
{\vergleichskette
{ P
}
{ \in }{ \partial M
}
{  }{
}
{    }{
}
{   }{
}
}
{}{}{}
ein Randpunkt. Zeige, dass für einen Tangentialvektor

\mavergleichskette
{\vergleichskette
{ v
}
{ \in }{ T_PM
}
{  }{
}
{    }{
}
{   }{
}
}
{}{}{}
folgende Eigenschaften äquivalent sind.
\aufzaehlungzwei {Es gibt einen stetig differenzierbaren Weg
\maabb {\gamma} { [- \epsilon, 0]} { M
} {}
mit

\mavergleichskette
{\vergleichskette
{ \gamma(0)
}
{ = }{ P
}
{  }{
}
{    }{
}
{   }{
}
}
{}{}{}
und

\mavergleichskette
{\vergleichskette
{ \gamma'(0)
}
{ = }{ v
}
{  }{
}
{    }{
}
{   }{
}
}
{}{}{.}
} { $v$ wird bei jeder Karte mit einem negativen Halbraum unter der Tangentialabbildung auf den rechten Halbraum abgebildet.
}

}
{} {}

\inputaufgabe
{}
{

Es sei $M$ eine
\definitionsverweis {differenzierbare Mannigfaltigkeit mit Rand}{}{}
und es sei
\maabbdisp {\gamma} {{]{-\epsilon}, \epsilon[}} {M
} {}
ein
\definitionsverweis {differenzierbarer Weg}{}{}
mit

\mavergleichskettedisp
{\vergleichskette
{ \gamma (0)
}
{ =} { P
}
{ \in} { \partial M
}
{ } {
}
{ } {
}
}
{}{}{.}
Zeige
\mathl{\gamma'(0) \in T_P \partial M}{.}

}
{} {}

\inputaufgabegibtloesung
{}
{

Es sei $M$ eine
\definitionsverweis {Mannigfaltigkeit mit Rand}{}{.}
Zeige, dass jede offene Teilmenge

\mavergleichskette
{\vergleichskette
{ N
}
{ \subseteq }{ M
}
{  }{
}
{    }{
}
{   }{
}
}
{}{}{}
ebenfalls eine Mannigfaltigkeit mit
\zusatzklammer {eventuell leerem} {} {}
Rand ist, und dass

\mavergleichskettedisp
{\vergleichskette
{ \partial N
}
{ =} { \partial M \cap N
}
{ } {
}
{ } {
}
{ } {
}
}
{}{}{}
gilt.

}
{} {}

\inputaufgabegibtloesung
{}
{

Es sei

\mavergleichskette
{\vergleichskette
{ H
}
{ \subseteq }{ \R^n
}
{  }{
}
{    }{
}
{   }{
}
}
{}{}{}
ein
\definitionsverweis {euklidischer Halbraum}{}{}
und

\mavergleichskette
{\vergleichskette
{ P
}
{ \in }{ H
}
{  }{
}
{    }{
}
{   }{
}
}
{}{}{.}
Es gebe eine in $\R^n$ offene Menge $V$ mit

\mavergleichskette
{\vergleichskette
{ P
}
{ \in }{ V
}
{ \subseteq }{ H
}
{    }{
}
{   }{
}
}
{}{}{.}
Zeige, dass
\mathl{P}{} kein Randpunkt von $H$ ist.

}
{} {}

\inputaufgabe
{}
{

Es seien
\mathl{U_1 \subseteq H_1 \subset \R^n}{}  und
\mathl{U_2 \subseteq H_2 \subset \R^n}{}
\definitionsverweis {offene Teilmengen}{}{}
in
\definitionsverweis {euklidischen Halbräumen}{}{}
\mathkor {} {H_1} {und} {H_2} {}
und es sei
\maabbdisp {\varphi} {U_1 } { U_2
} {}
ein
\definitionsverweis {Diffeomorphismus}{}{.}
Zeige, dass es zu jedem Punkt
\mathl{P\in U_1}{} offene Umgebungen
\mathl{P\in V_1 \subset \R^n}{} und
\mathl{\varphi(P) \in V_2 \subset \R^n}{} und eine diffeomorphe Fortsetzung
\maabbdisp {\tilde{\varphi}} {V_1 } { V_2
} {}
gibt.

}
{} {}

\inputaufgabe
{}
{

Die abgeschlossene Kreisscheibe
\mathl{B  \left(  0 , 1 \right)}{} trage die
\definitionsverweis {Standardorientierung}{}{}
des $\R^2$. Läuft die durch die
\definitionsverweis {äußere Normale}{}{}
festgelegte Orientierung auf dem Rand
\zusatzklammer {also auf dem Einheitskreis} {} {} mit dem oder gegen den Uhrzeigersinn?

}
{} {}

  \zwischenueberschrift{Aufgaben zum Abgeben}

\bild{ \begin{center}
\includegraphics[width=5.5cm]{ \bildeinlesungsvg {Not-star-shaped} {svg} }
\end{center}
\bildtext {} }

\bildlizenz { Not-star-shaped.svg } {} {Oleg Alexandrov} {Commons} {PD} {}

\inputaufgabe
{4}
{

Man gebe für den Kreisring

\mavergleichskettedisp
{\vergleichskette
{M
}
{ =} { \left\{  x \in \R^2  \mid   1 \leq \Vert { x } \Vert \leq 2         \right\}
}
{ } {
}
{ } {
}
{ } {
}
}
{}{}{}
explizit
\definitionsverweis {Karten}{}{}
an, die zeigen, dass $M$ eine
\definitionsverweis {Mannigfaltigkeit mit Rand}{}{}
ist.

}
{} {}

\inputaufgabe
{6}
{

Zeige, dass die
\definitionsverweis {Halbebene}{}{} $\R_{\geq 0} \times \R$ und der
\definitionsverweis {Quadrant}{}{}
\mathl{\R_{\geq 0} \times \R_{\geq 0}}{} nicht
\definitionsverweis {diffeomorph}{}{}
sind.

}
{(Was ist hierbei der geeignete Diffeomorphiebegriff?)} {}

\inputaufgabe
{6}
{

Es sei
\mathl{M= B  \left(  0 , 1 \right) \setminus \{ (0,1), (0,-1)\} \subseteq \R^2}{,} also die
\definitionsverweis {abgeschlossene Kreisscheibe}{}{,} aus der man zwei Randpunkte herausgenommen hat. Es sei
\mathl{N={]{-1},1[} \times  [-1,1]}{} das Produkt eines offenen und eines abgeschlossenen Intervalls. Zeige, dass
\mathkor {} {M} {und} {N} {}
\definitionsverweis {diffeomorphe}{}{}
\definitionsverweis {Mannigfaltigkeiten mit Rand}{}{} sind.

}
{} {}

\inputaufgabe
{4}
{

Es seien
\mathl{V_1,V_2 \subseteq \R^n}{}
\definitionsverweis {offene Teilmengen}{}{}
und es sei
\maabbdisp {\varphi} {V_1 } { V_2
} {}
ein
\definitionsverweis {Diffeomorphismus}{}{,}
der eine
\definitionsverweis {Homöomorphie}{}{}
zwischen
\mathkor {} {V_1 \cap H} {und} {V_2 \cap H} {}
induziert und damit auch zwischen
\mathkor {} {V_1 \cap \partial H} {und} {V_2 \cap \partial H} {}
\zusatzklammer {$H$ bezeichnet den
\definitionsverweis {Halbraum}{}{}
und $\partial H$ seinen Rand} {} {.}
Zeige, dass die Einschränkung auf den Rand ebenfalls ein
\definitionsverweis {Diffeomorphismus}{}{}
ist.

}
{} {}

\inputaufgabe
{4}
{

Die abgeschlossene Einheitskugel
\mathl{B  \left(  0 , 1 \right)}{} trage die
\definitionsverweis {Standardorientierung}{}{} des $\R^3$. Bestimme, ob die beiden Tangentenvektoren
\mathkor {} {(2,1,0)} {und} {(3,-1,0)} {}
am Nordpol
\mathl{(0,0,1)}{} die durch die
\definitionsverweis {äußere Normale}{}{} induzierte Orientierung auf dem Rand
\zusatzklammer {also auf der Einheitssphäre} {} {} repräsentieren oder nicht?

}
{} {}

 [[Kategorie:Latexseite]]
